\documentclass[11pt]{article}
\usepackage[margin=1.15in]{geometry}
\usepackage{amsmath,amssymb,amsthm,mathtools}
\usepackage{microtype}

\theoremstyle{plain}
\newtheorem{theorem}{Theorem}
\newtheorem{proposition}[theorem]{Proposition}
\theoremstyle{definition}
\newtheorem{claim}[theorem]{Printed statement}
\theoremstyle{remark}
\newtheorem*{remark}{Remark}

\DeclareMathOperator{\Var}{Var}
\DeclareMathOperator{\Cov}{Cov}
\DeclareMathOperator{\E}{\mathbb{E}}
\newcommand{\R}{\mathbb{R}}
\newcommand{\X}{\mathcal{X}}
\newcommand{\Pt}{P_\theta}
\newcommand{\score}{\dot\ell_\theta}

\title{Two defective statements in \textsc{StatLean/PointEstimation}\\
\large The final two open goals of Batch 11 are not formalization gaps}
\author{StatLean formalization notes}
\date{26 July 2026}

\begin{document}
\maketitle

\begin{abstract}
The \textsc{PointEstimation} area (Lehmann--Casella, \emph{Theory of Point Estimation},
2nd ed.) is closed except for two goals. Neither is a missing proof: in both cases the
statement as transcribed from the source is \emph{defective}, and we verified this rather
than assumed it. The first, the converse half of the attainment characterisation
(TPE2~\S2.5, Thm.~5.12), is \textbf{false as written}; we give an explicit counterexample
and isolate the missing hypothesis. The second, the nonexistence of a best linear unbiased
estimator under a random design (Shaffer's phenomenon, TPE2~\S3.4), is \textbf{vacuously
unprovable as written} because its hypothesis is strictly too weak; we give a
counterexample to the printed form and state the corrected hypothesis. Closing either
goal is therefore a decision about what to state, not an exercise in proving what is
stated.
\end{abstract}

\section{Attainment of the information bound forces an exponential family}

\subsection{The printed statement}

Let $\{\Pt\}_{\theta\in\R}$ be a family of probability measures on a measurable space
$\X$ with densities $p_\theta$ with respect to a common dominating measure $\mu$, all with
the same support. Write
\[
  \score(x) \;=\; \frac{\partial}{\partial\theta}\log p_\theta(x),
  \qquad
  I(\theta) \;=\; \E_\theta\!\left[\score(X)^2\right]
\]
for the score and the Fisher information, and let $\delta:\X\to\R$ be a square-integrable
statistic estimating
\[
  g(\theta) \;=\; \E_\theta\,\delta(X), \qquad g'(\theta) \;=\; \frac{d}{d\theta}\,g(\theta).
\]
The information inequality states $\Var_\theta(\delta)\ \ge\ g'(\theta)^2/I(\theta)$.

\begin{claim}[TPE2 \S2.5, Thm.~5.12, converse half]\label{claim:attain}
Suppose $I(\theta)>0$ and $\E_\theta\,\score = 0$ for every $\theta$, and suppose the bound
is attained at \emph{every} parameter value:
\begin{equation}\label{eq:attain}
  \Var_\theta(\delta) \;=\; \frac{g'(\theta)^2}{I(\theta)}
  \qquad\text{for all }\theta .
\end{equation}
Then $\{\Pt\}$ is a one-parameter exponential family with $\delta$ as its natural
statistic: there exist $\eta\in C^1(\R)$, $B:\R\to\R$ and a measurable $h\ge 0$ with
\begin{equation}\label{eq:expform}
  p_\theta(x) \;=\; \exp\!\bigl(\eta(\theta)\,\delta(x) - B(\theta)\bigr)\,h(x)
  \qquad \text{for } \mu\text{-a.e. } x,\ \text{every }\theta .
\end{equation}
\end{claim}

(Our transcription additionally carries regularity --- joint measurability of
$(\theta,x)\mapsto p_\theta(x)$, $C^1$ dependence on $\theta$ for a.e.\ $x$, and continuity
of $I$ and of $g'$ --- which only strengthens the hypotheses and so is irrelevant to the
refutation below.)

\subsection{The counterexample}

\begin{proposition}\label{prop:attain-false}
Claim~\ref{claim:attain} is false. There is a family satisfying every hypothesis, including
\eqref{eq:attain} at every $\theta$, for which no representation \eqref{eq:expform} exists.
\end{proposition}

\begin{proof}
Take $\X=\R$, $\mu$ Lebesgue measure, the standard Gaussian location family
\[
  p_\theta(x) \;=\; \varphi(x-\theta) \;=\; \frac{1}{\sqrt{2\pi}}\,e^{-(x-\theta)^2/2},
\]
and the estimator
\[
  \delta \;\equiv\; 0 .
\]

\emph{All hypotheses hold.} The densities are strictly positive, so the common-support
condition is satisfied, and $(\theta,x)\mapsto\varphi(x-\theta)$ is jointly measurable and
$C^\infty$ in $\theta$. The score is
\[
  \score(x) \;=\; \frac{\partial}{\partial\theta}\left(-\tfrac{1}{2}(x-\theta)^2\right)
  \;=\; x-\theta,
\]
whence $\E_\theta\,\score = 0$ and $I(\theta)=\E_\theta (X-\theta)^2 = 1 > 0$; both are
continuous in $\theta$. The statistic $\delta\equiv 0$ is square-integrable, and
\[
  g(\theta) \;=\; \E_\theta\,\delta \;=\; 0 \quad\text{for all }\theta,
  \qquad\text{so}\qquad g'(\theta)\;=\;0,
\]
which is continuous. Finally $\Var_\theta(\delta)=0$, so the attainment identity
\eqref{eq:attain} reads
\[
  0 \;=\; \frac{0^2}{1} \;=\; 0 ,
\]
and holds at every $\theta$.

\emph{The conclusion fails.} Suppose \eqref{eq:expform} held. Since $\delta\equiv 0$ the
exponential factor does not depend on $x$, so for every $\theta$ there is a constant
$c(\theta) = e^{-B(\theta)}$ with
\[
  \varphi(x-\theta) \;=\; c(\theta)\,h(x) \qquad\text{for Lebesgue-a.e.\ }x .
\]
Comparing $\theta=1$ with $\theta=0$ and dividing (legitimate, as $\varphi>0$ forces
$h>0$ a.e.) gives
\[
  \frac{\varphi(x-1)}{\varphi(x)} \;=\; \frac{c(1)}{c(0)}
  \qquad\text{for a.e.\ }x .
\]
But the left-hand side is
\[
  \frac{\varphi(x-1)}{\varphi(x)}
  \;=\; \exp\!\left(-\tfrac{(x-1)^2}{2}+\tfrac{x^2}{2}\right)
  \;=\; e^{\,x-\frac12},
\]
which is strictly increasing and therefore not a.e.\ equal to any constant. Contradiction.

A finite counterexample is available with the same $\delta\equiv 0$: take
$\X=\{0,1\}$ with the Bernoulli family $p_\theta$, where the same computation applies.
\end{proof}

\subsection{Diagnosis: the attainment hypothesis is vacuous at stationary points}

The information inequality is Cauchy--Schwarz applied to $\delta-\E_\theta\delta$ and
$\score$:
\begin{equation}\label{eq:cs}
  \Cov_\theta(\delta,\score)^2 \;\le\; \Var_\theta(\delta)\; I(\theta),
\end{equation}
combined with the identity $\Cov_\theta(\delta,\score)=g'(\theta)$, which is where
differentiation under the integral sign enters. Equality in \eqref{eq:cs} holds if and only
if $\delta-\E_\theta\delta$ and $\score$ are almost surely proportional --- and it is
\emph{that} proportionality, not the numerical identity \eqref{eq:attain}, which drives the
classical proof: it makes $\score$ an affine function of $\delta$, and integrating in
$\theta$ produces \eqref{eq:expform}.

The defect is that \eqref{eq:attain} degenerates precisely when
\[
  g'(\theta)=0 \quad\text{and}\quad \Var_\theta(\delta)=0 .
\]
Then \eqref{eq:attain} reads $0=0$; the Cauchy--Schwarz equality case is attained
trivially, with $\delta$ almost surely constant, and it constrains the score \emph{not at
all}. The hypothesis is satisfied by every family whatsoever, paired with a constant
$\delta$ --- which is exactly what the counterexample exploits.

The classical statement is about an estimator of an estimand that genuinely moves. Lehmann
and Casella state Thm.~5.12 alongside the companion identities
\[
  \delta \;=\; \frac{g'(\theta)}{I(\theta)}\,\score + g(\theta),
  \qquad
  I(\theta) \;=\; \eta'(\theta)\,g'(\theta),
\]
the second of which is meaningless unless $g'(\theta)\neq 0$; the nondegeneracy is carried
implicitly by the surrounding text rather than by the displayed hypotheses. Our
transcription kept the displayed hypotheses and lost the implicit one.

\subsection{The repair}

Add the nondegeneracy hypothesis
\[
  g'(\theta) \;\neq\; 0 \qquad\text{for every }\theta .
\]
With it, the equality case of \eqref{eq:cs} is no longer vacuous and the classical route
runs. Two of its three steps are already formalized in the area:

\begin{enumerate}
\item \emph{From attainment to Cauchy--Schwarz equality.} This requires
  $\Cov_\theta(\delta,\score)=g'(\theta)$, i.e.\ the interchange
  $\frac{d}{d\theta}\int \delta\,p_\theta \,d\mu = \int \delta\,\partial_\theta p_\theta\,d\mu$.
  This is \emph{not} implied by the present hypotheses, which give only the value of the
  derivative of $\theta\mapsto\E_\theta\delta$, not the interchange; without it one obtains
  only $\Cov^2\le g'^2$. A domination hypothesis must be supplied.
\item \emph{From equality to proportionality.} Formalized: given the interchange and
  $g'(\theta)\ne 0$ one has $\Var_\theta\bigl(\delta - c\,\score\bigr)=0$ with
  $c=g'(\theta)/I(\theta)$, hence
  $\score = \frac{I(\theta)}{g'(\theta)}\bigl(\delta-\E_\theta\delta\bigr)$ almost surely.
\item \emph{Integration in the parameter, uniformly in $x$.} For a.e.\ $x$ the map
  $\theta\mapsto\log p_\theta(x)$ is $C^1$ with derivative $\eta'(\theta)\delta(x)-B'(\theta)$;
  the fundamental theorem of calculus then yields \eqref{eq:expform}. No library lemma
  packages this uniform-in-$x$ integration at the current Mathlib pin.
\end{enumerate}

Consequently the goal is left open deliberately, with the falsity of the printed form
recorded at the statement. Repairing it is a signature change, and a decision for the
library, not a proof obligation.

\section{Nonexistence of a BLUE under a random design}

\subsection{The setting and the printed statement}

Let the design be a random element $a$ of $(\R^{s})^{n}$, written as a matrix
$A=A(a)\in\R^{s\times n}$, with law $Q$ free of the parameter. Given the design, the
response $Y\in\R^{n}$ has
\[
  \E\bigl[Y \mid a\bigr] \;=\; A^{\!\top}\theta,
  \qquad
  \Cov\bigl[Y \mid a\bigr] \;=\; \sigma^{2} I_n ,
\]
with unknown parameter $p=(\theta,\sigma^{2})\in\R^{s}\times(0,\infty)$; the joint law is
$P_p = Q\otimes K_p$. A \emph{linear} estimator is one of the form
\[
  \delta(a,y) \;=\; \langle c(a),\,y\rangle ,
\]
its coefficient vector allowed to depend on the observed design. Fix $c\in\R^{s}$,
$c\ne 0$, and estimate the functional $c^{\!\top}\theta$. A \emph{best linear unbiased
estimator} (BLUE) is a linear unbiased square-integrable $\delta$ whose variance is minimal,
at \emph{every} parameter value, among all linear unbiased square-integrable competitors.

\begin{claim}[TPE2 \S3.4, Shaffer's phenomenon]\label{claim:blue}
If the design law $Q$ is not a point mass, then no best linear unbiased estimator of
$c^{\!\top}\theta$ exists.
\end{claim}

\subsection{The variance decomposition}

Write $A=A(a)$ and put
\[
  v(a) \;:=\; A(a)\,c(a) \in \R^{s},
  \qquad
  i(a) \;:=\; c^{\!\top}\bigl(A(a)A(a)^{\!\top}\bigr)^{-1}c ,
\]
so that $i(a)$ is the \emph{design information for the functional $c$} (the reciprocal of
the conditional least-squares variance, up to $\sigma^2$). Conditionally on the design,
$\E[\delta\mid a]=\langle c(a), A^{\!\top}\theta\rangle = \theta^{\!\top}v(a)$, so
unbiasedness for $c^{\!\top}\theta$ at every $\theta$ is
\begin{equation}\label{eq:unb}
  \E_Q\,v \;=\; c .
\end{equation}
By the law of total variance,
\begin{equation}\label{eq:ltv}
  \Var(\delta)
  \;=\; \sigma^{2}\,\E_Q\bigl\|c(a)\bigr\|^{2}
      \;+\; \Var_Q\bigl(\theta^{\!\top}v(a)\bigr) .
\end{equation}
For fixed $v(a)$, minimising $\|c(a)\|^{2}$ subject to $A(a)c(a)=v(a)$ gives
$\|c(a)\|^{2}\ \ge\ v(a)^{\!\top}(AA^{\!\top})^{-1}v(a)$, with equality for the conditional
least-squares choice. Hence, writing $v=\lambda c$ in the scalar case $s=1$,
\begin{equation}\label{eq:lb}
  \Var(\delta)
  \;\ge\; \sigma^{2}\,\E_Q\bigl[\lambda^{2}\, i\bigr]
      \;+\; \theta^{2}\,\Var_Q(\lambda),
  \qquad \E_Q\lambda = 1 .
\end{equation}

Shaffer's phenomenon is the observation that the two terms of \eqref{eq:lb} are minimised by
\emph{different} weightings:
\[
  \min \E_Q[\lambda^{2} i] \text{ subject to } \E_Q\lambda=1
  \ \Longrightarrow\ \lambda \propto 1/i,
  \qquad
  \min \Var_Q(\lambda) \text{ subject to } \E_Q\lambda=1
  \ \Longrightarrow\ \lambda \equiv 1 .
\]
These two prescriptions coincide if and only if $i$ is $Q$-almost surely constant. When $i$
genuinely varies, no single linear unbiased estimator minimises the variance at all
$(\theta,\sigma^{2})$ simultaneously, and no BLUE exists.

\subsection{The counterexample to the printed hypothesis}

\begin{proposition}\label{prop:blue-false}
Claim~\ref{claim:blue} is false: ``$Q$ is not a point mass'' does not imply nonexistence.
\end{proposition}

\begin{proof}
Take $s=1$, $n=2$, $c=1$, fix $a\in\R^{2}\setminus\{0\}$ and let
\[
  Q \;=\; \tfrac12\,\delta_{a} \;+\; \tfrac12\,\delta_{-a},
  \qquad
  K_p(a) \;=\; \mathcal{N}\bigl(\theta a,\ \sigma^{2}I_2\bigr).
\]
Then $Q$ is not a point mass, so the printed hypothesis holds.

The two designs $a$ and $-a$ have the same norm, hence the same information
\[
  i(a) \;=\; i(-a) \;=\; \frac{1}{\|a\|^{2}} .
\]
So the lower bound \eqref{eq:lb} becomes
\[
  \Var(\delta)
  \;\ge\; \frac{\sigma^{2}}{2\|a\|^{2}}\Bigl(\lambda(a)^{2}+\lambda(-a)^{2}\Bigr)
      \;+\; \theta^{2}\,\Var_Q(\lambda),
  \qquad
  \frac{\lambda(a)+\lambda(-a)}{2}=1 .
\]
Subject to that constraint, the first term is minimised at
$\lambda(a)=\lambda(-a)=1$ (the constrained minimum of a sum of squares with equal weights),
and the second term $\Var_Q(\lambda)$ is minimised, with value $0$, at the same point.
\emph{Both terms are therefore minimised simultaneously}, and the conditional least-squares
estimator --- $\lambda\equiv 1$, i.e.\ $c(a)=a/\|a\|^{2}$, that is
$\delta(a,y)=\langle a,y\rangle/\|a\|^{2}$ --- attains the minimum variance at every
$(\theta,\sigma^{2})$. It is linear, unbiased and square-integrable, hence a BLUE. The
printed negation fails.
\end{proof}

\begin{remark}
The example is not pathological: it is the generic situation for any design law that is
non-degenerate but \emph{informationally homogeneous}. The printed hypothesis measures the
wrong thing --- variability of the design --- where the argument needs variability of the
information the design carries about $c^{\!\top}\theta$.
\end{remark}

\subsection{The repair}

Replace the printed hypothesis by the substantive one that the argument actually uses:
\[
  \boxed{\ \text{the design information } a\ \mapsto\ c^{\!\top}\bigl(A(a)A(a)^{\!\top}\bigr)^{-1}c
  \ \text{ is not } Q\text{-almost surely constant.}\ }
\]
This is strictly stronger than the printed condition --- a point mass has constant
information, so the new hypothesis implies the old --- and it is exactly the dichotomy
exposed by \eqref{eq:lb}. A regularity companion is also required, namely that almost every
realised design have full row rank $s$, so that $(AA^{\!\top})^{-1}$ is the genuine Gram
inverse.

With the amended hypothesis the statement is believed true, and the classical route is:
if $\delta$ is a BLUE and $d$ is linear and unbiased for $0$ at every parameter, then
$t\mapsto\Var(\delta+t\,d)$ is a nonnegative quadratic minimised at $t=0$, so
$\Cov(\delta,d)=0$ at every parameter. Choosing
$d(a,y)=\psi(a)\,\langle A^{\!\top}(AA^{\!\top})^{-1}c,\ y\rangle$ for bounded $\psi$ with
$\E_Q\psi=0$ and splitting the covariance into its $\sigma^{2}$- and $\theta^{2}$-parts
forces first $v=c$ almost surely, and then $i$ to be almost surely constant --- contradicting
the amended hypothesis. Discharging this in the library additionally requires a measurability
conjunct on the coefficient map $a\mapsto c(a)$, which the definition of a linear estimator
does not currently carry.

\section{Summary}

\begin{center}
\begin{tabular}{@{}p{0.30\textwidth}p{0.62\textwidth}@{}}
\hline
\textbf{Goal} & \textbf{Status} \\
\hline
Attainment $\Rightarrow$ exponential family (\S2.5, Thm.~5.12) &
\textbf{False as written.} The attainment identity is vacuous where $g'(\theta)=0$ and
$\Var_\theta(\delta)=0$; the Gaussian location family with $\delta\equiv 0$ satisfies every
hypothesis and is not exponential in $\delta$. Repair: assume $g'\ne 0$ throughout. One
analytic ingredient (uniform-in-$x$ parametric integration) would still be needed. \\[1.2ex]
Nonexistence of a BLUE, random design (\S3.4) &
\textbf{Hypothesis too weak.} ``$Q$ not a point mass'' admits informationally homogeneous
designs such as $\tfrac12\delta_a+\tfrac12\delta_{-a}$, for which conditional least squares
\emph{is} a BLUE. Repair: assume the design information $c^{\!\top}(AA^{\!\top})^{-1}c$ is
not a.s.\ constant. \\
\hline
\end{tabular}
\end{center}

\noindent
Every other result of the area is proved and the area builds green. In both cases above the
remaining work is to decide the corrected statement --- an editorial act on the
formalization's contract with the source --- after which the mathematics is routine in the
first case and classical in the second.

\end{document}
